\section{Troubleshooting}

\subsection{Common Issues}

\subsubsection{Import Errors}
\begin{lstlisting}[language=Python, caption={Path Setup}]
# Ensure correct path setup
import sys
sys.path.insert(0, '/path/to/BioNetFlux/code')
\end{lstlisting}

\subsubsection{Geometry Validation}
\begin{lstlisting}[language=Python, caption={Geometry Debugging}]
# Check geometry before problem creation
geometry = DomainGeometry("test")
# ... add domains ...
print(geometry.summary())  # Verify domain layout
print(geometry.get_bounding_box())  # Check coordinates
\end{lstlisting}

\subsubsection{Constraint Setup}
\begin{lstlisting}[language=Python, caption={Constraint Verification}]
# Verify constraint mapping
constraint_manager.map_to_discretizations(discretizations)
print(f"Total constraints: {constraint_manager.n_multipliers}")
\end{lstlisting}

\subsubsection{Solution Convergence}
\begin{lstlisting}[language=Python, caption={Convergence Monitoring}]
# Monitor Newton iteration
newton_tolerance = 1e-10
max_newton_iterations = 20

# Check residual norms during iteration
if residual_norm > newton_tolerance:
    print(f"Convergence issue: residual = {residual_norm:.2e}")
\end{lstlisting}

\subsection{Performance Optimization}

\begin{enumerate}
    \item \textbf{Mesh Resolution}: Balance accuracy vs. computational cost
    \item \textbf{Time Step Size}: Use adaptive time stepping for stability
    \item \textbf{Newton Tolerance}: Adjust based on problem requirements
    \item \textbf{Domain Decomposition}: Optimize domain sizes for load balancing
\end{enumerate}

\subsection{Debugging Tips}

\begin{enumerate}
    \item \textbf{Visualization}: Use all three plot types to understand solution behavior
    \item \textbf{Parameter Validation}: Check physical parameter ranges
    \item \textbf{Constraint Verification}: Ensure proper interface connectivity
    \item \textbf{Solution Monitoring}: Track solution norms and residuals
\end{enumerate}
